%% ============================================================
%%  reproducibilidad.tex  —  Sección 1.7
%% ============================================================

\subsection{Reproducibilidad}

Todos los componentes del proyecto están diseñados para ser totalmente
reproducibles. A continuación se detallan los pasos necesarios para regenerar
los datos, gráficas e informe desde cero.

\subsubsection*{Estructura del repositorio}

\begin{verbatim}
Proyecto_Avanzados/
(graph.h, dijkstra.h, main.cpp)
|-- paper_1_dijkstra/src/
(bellman_ford.h, main.cpp)
|-- paper_2_bellman_ford/src/
(thorup.h, main.cpp)
|-- paper_3_thorup/src/
(det_mlogn.h, main.cpp)
|-- paper_4_henzinger/src/
|-- aporte_propio/
(todos los .h unificados + benchmarks)
|   |-- src/
|   |   |-- graph.h
|   |   |-- dijkstra.h
|   |   |-- bellman_ford.h
|   |   |-- thorup.h
|   |   |-- det_mlogn.h
        (exps. 1-3)
|   |   |-- main_benchmark.cpp
        (exp. 4 - aporte propio)
|   |   |-- density_experiment.cpp
        (generador de gráficas)
|   |   `-- plot_density.py
|   |-- resultados/
|   |   |-- densidad_raw.csv
|   |   `-- densidad_promedio.csv
(4 PNG generados por plot_density.py)
|   |-- graficas/
    (resultados exps. 1-3)
|   |-- results.csv
    (análisis textual exps. 1-3)
|   |-- analisis.md
    (análisis exp. 4)
|   `-- analisis_densidad.md
\end{verbatim}

\subsubsection*{Paso 1: Compilar y ejecutar los benchmarks individuales}

Desde cada carpeta \texttt{src/} individual:

\begin{lstlisting}[language=bash, caption={Compilación individual de cada algoritmo.}, captionpos=b]
# Desde paper_1_dijkstra/src/
g++ -O2 -std=c++17 -Wall -o dijkstra_bench main.cpp
./dijkstra_bench

# Desde paper_2_bellman_ford/src/
g++ -O2 -std=c++17 -Wall -o bellman_ford_bench main.cpp
./bellman_ford_bench

# Desde paper_3_thorup/src/
g++ -O2 -std=c++17 -Wall -o thorup_bench main.cpp
./thorup_bench

# Desde paper_4_henzinger/src/
g++ -O2 -std=c++17 -Wall -o dmmsy_bench main.cpp
./dmmsy_bench
\end{lstlisting}

\subsubsection*{Paso 2: Compilar y ejecutar el benchmark unificado (Exps.\ 1–3)}
~\\
\begin{lstlisting}[language=bash,
  caption={Benchmark unificado — genera \texttt{results.csv}.}, captionpos=b]
# Desde aporte_propio/src/
g++ -O2 -std=c++17 -Wall -o sssp_benchmark main_benchmark.cpp
./sssp_benchmark
# Salida: aporte_propio/results.csv
\end{lstlisting}
~\\
\subsubsection*{Paso 3: Compilar y ejecutar el experimento de densidad (Exp.\ 4)}
~\\
\begin{lstlisting}[language=bash,
  caption={Experimento de densidad — genera dos CSV en \texttt{resultados/}.}, captionpos=b]
# Desde aporte_propio/src/
g++ -O2 -std=c++17 -Wall -o density_experiment density_experiment.cpp
./density_experiment
# Salida:
#   aporte_propio/resultados/densidad_raw.csv
#   aporte_propio/resultados/densidad_promedio.csv
\end{lstlisting}

El programa imprime progreso en consola para cada configuración
$(n, \text{densidad})$ e indica las rutas de los archivos generados al
finalizar.
~\\
\subsubsection*{Paso 4: Generar las gráficas}
~\\
\begin{lstlisting}[language=bash,
  caption={Generación de gráficas PNG.}, captionpos=b]
# Requiere: pip install matplotlib pandas seaborn
# Desde aporte_propio/src/
python plot_density.py
# Salida: aporte_propio/graficas/*.png (4 archivos)
\end{lstlisting}
~\\

\subsubsection*{Factores que pueden afectar la reproducibilidad}

\begin{itemize}
  \item \textbf{Semillas:} Los experimentos 1–3 usan la semilla fija $42$.
    El Experimento~4 usa las semillas $\{42, 123, 999\}$.
    Cambiar la semilla altera los grafos generados y por tanto los tiempos,
    aunque el número de aristas y la densidad se mantienen en el rango
    objetivo.
  \item \textbf{Versión del compilador:} El uso de \texttt{-O2} activa
    optimizaciones dependientes del compilador. Diferentes versiones de
    \texttt{g++} pueden generar código con tiempos distintos.
  \item \textbf{Carga del sistema:} Los tiempos en microsegundos son
    sensibles a la carga del sistema operativo durante la medición. Se
    recomienda ejecutar los benchmarks con carga mínima de procesos en
    segundo plano.
  \item \textbf{Arquitectura del procesador:} Las implementaciones no usan
    instrucciones SIMD ni paralelismo; sin embargo, la jerarquía de caché
    del procesador influye en los tiempos reales.
\end{itemize}
